\chapter{Error handling}
\label{ch:fel}

{\large\itshape Five kinds of error, one signal, and a node that switches itself off.\par}

\begin{chapterbox}{In this chapter}
\begin{itemize}
  \item The five errors a CAN node can detect, and who detects what.
  \item The error frame: how a node aborts a transfer in progress.
  \item The error counters, and the three states a node can be in.
  \item Bus-off: when a node disconnects itself, and why that is the right behaviour.
\end{itemize}
\end{chapterbox}

\section{Five kinds of error}

CAN detects errors in five ways. Three of them we have already covered under other names:

\begin{booktable}{@{}lLl@{}}
\thead{Error} & \thead{What it is} & \thead{Who sees it} \\ \midrule
Bit error & A transmitting node reads back something other than what it sent & The transmitter \\
Stuff error & Six equal bits in a row where the rule of five demands an edge & Everyone \\
CRC error & The received checksum does not match & Receivers \\
Form error & A field that must be recessive is not & Everyone \\
Acknowledgement error & Nobody pulled down the ACK slot & The transmitter \\
\end{booktable}

Together they give a residual error probability usually quoted at around $10^{-11}$ of the errors
that occur --- roughly one undetected error per thousand years on a heavily loaded bus. That is not
luck but the result of five independent mechanisms overlapping.

\subsection{Bit errors, and the exceptions}

A transmitting node constantly compares what it sends with what is on the bus. If they differ it is
a bit error --- \emph{except} in two places where a difference is perfectly normal:

\begin{itemize}
  \item \textbf{During the arbitration field.} Recessive out and dominant in means lost
    arbitration, not an error (\kapref{ch:arbitrering}).
  \item \textbf{In the ACK slot.} The transmitter sends recessive and \emph{hopes} for dominant
    (\kapref{ch:frameformat}).
\end{itemize}

Telling those two cases apart from a real bit error is one of the things that make a CAN controller
more involved than it looks.

\section{The error frame}

When a node detects an error it does not say so afterwards. It aborts the frame \textbf{while it is
in progress}, by sending an \textbf{error frame}: six dominant bits in a row.

Six dominant bits are impossible in a correct frame --- the rule of five sees to it that there are
never more than five --- so the pattern cannot be mistaken for data. Every other node on the bus
sees it, understands that the frame is broken, and discards it.

And since every node sees the error, several of them will send their own error flag in turn. The
result is a run of dominant bits, followed by a delimiter of eight recessive ones.

Two things follow:

\textbf{Every node discards the frame, not just the one that detected the error.} The bus is kept
consistent: either everyone receives the frame, or nobody does. That is a stronger guarantee than a
home-grown protocol normally attempts.

\textbf{The transmitter resends automatically.} The controller does it by itself, without the
software being involved. A driver that sends a frame and learns that it went well may therefore have
had it resent several times without noticing.

\section{The error counters}

Here comes the part of CAN that is least obvious and most carefully thought out.

A node that runs into errors should speak up. A node that is \emph{broken} should stop speaking up
--- otherwise it ruins the bus for everyone else by sending error frames non-stop.

CAN tells them apart with two counters in every node:

\begin{itemize}
  \item \textbf{TEC} --- transmit error counter, errors while the node was sending.
  \item \textbf{REC} --- receive error counter, errors while it was receiving.
\end{itemize}

The rules in brief: an error increases the counter by 8 (transmitter side) or 1 (receiver side),
and a successful transfer decreases it by 1.

\textbf{The asymmetry is the whole point.} A node that really is broken is involved in every error
and climbs fast. A node that just happens to be on a noisy bus climbs slowly and has time to work it
off. So the counters do not measure ``how many errors have happened'' but ``how likely is it that
the error is mine''.

\subsection{The three states}

\bookfigure{error_states}{A node's three error states and the transitions between them.}{fig:errorstates}

\begin{booktable}{@{}llL@{}}
\thead{State} & \thead{Condition} & \thead{Behaviour} \\ \midrule
Error active & TEC and REC $\leq$ 127 & Normal operation; sends \emph{active} error flags (dominant) \\
Error passive & Either counter $>$ 127 & Sends \emph{passive} error flags (recessive) and waits extra between transmissions \\
Bus off & TEC $>$ 255 & Sends nothing at all \\
\end{booktable}

The transition to \emph{error passive} is the interesting one. A passive error flag is recessive,
which means it does not disturb anyone else. The node may still take part, it may still report,
but it can no longer pull the bus down on everyone's behalf.

\subsection{Bus off}

With TEC above 255 the node disconnects itself. It does not send, does not acknowledge, does not
disturb.

It is the protocol's last line of defence, and it is worth dwelling on how unusual the behaviour
is: a node that concludes that it is itself most likely the broken part, and that draws that
conclusion without anyone telling it.

A node in bus off can be brought back --- typically after seeing the bus idle a certain number of
times, or after the software reinitialises the controller. That this recovery exists is deliberate
too: the fault may have been a loose connector that is now seated.

\begin{aside}
\note[For the driver writer] Bus off is the one state the software cannot ignore. Everything else
the controller handles by itself --- retransmissions, error frames, counters --- but a node in bus
off never comes back on its own if the design requires a reinitialisation. It is one of the few
places where error handling reaches all the way up into the application.
\end{aside}

\section{What the software sees}

Notice how little of all this reaches the software in a complete controller:

\begin{itemize}
  \item Errors are detected in hardware.
  \item Error frames are sent in hardware.
  \item Retransmission happens in hardware.
  \item The counters are updated in hardware.
\end{itemize}

The software gets a status bit, perhaps an error code, and the ability to read the counters. That
is an enormous amount of work done on its behalf --- compare with ACK, timeout, retry and duplicate
handling in a home-grown protocol, which are the same problems solved in software, worse.

\begin{simplification}
\note[Simplification] The simplified controller in \kapref{ch:kursen} implements \emph{nothing} of
this chapter except stuff error detection and the CRC check. It sends no error frames, has no error
counters, knows nothing of error passive or bus off, and never resends automatically. It only
signals that something went wrong, with a single bit, and leaves the rest to the software.
\end{simplification}

\section*{Exercises}

\exercise{Who sees the error?}{Reasoning}
\label{ex:6:1}
For each of the following: what kind of error is it, and who detects it?
\begin{enumerate}[a)]
  \item A disturbance flips a bit in the data field.
  \item A disturbance flips a bit in a field that must be recessive.
  \item The transmitter is alone on the bus.
  \item Seven equal bits in a row go by in the arbitration field.
\end{enumerate}

\exercise{The counters}{Calculation}
\label{ex:6:2}
A node starts at TEC 0 and runs into eight transmit errors in a row, followed by twenty successful
transmissions.
\begin{enumerate}[a)]
  \item What is TEC after the errors?
  \item After the successful transmissions?
  \item Which state is the node in at each of those points?
\end{enumerate}

\exercise{The asymmetry}{Reasoning}
\label{ex:6:3}
An error increases the transmitter's counter by 8 but the receiver's by 1.
\begin{enumerate}[a)]
  \item Why are the increases different?
  \item A node with a broken transmitter and a node on a noisy bus look alike at first. What tells
    them apart after a while?
\end{enumerate}

\exercise{Bus off}{Reasoning}
\label{ex:6:4}
\begin{enumerate}[a)]
  \item Why does a node disconnect itself instead of carrying on reporting errors?
  \item What would have happened to the bus if it did not?
  \item Why is there a way back from bus off?
\end{enumerate}
